\documentclass[11pt]{article}
\usepackage[a4paper,margin=25mm]{geometry}
\usepackage{amsmath,amssymb,mathtools}
\usepackage{fontspec,unicode-math}
\setmainfont{Latin Modern Roman}
\setsansfont{Latin Modern Sans}
\setmonofont{Latin Modern Mono}
\setmathfont{Latin Modern Math}
\usepackage{xurl,hyperref}
\hypersetup{colorlinks=true,urlcolor=blue,linkcolor=blue}
\newcommand{\sref}[2]{\hyperlink{source:#1:#2}{[S#2]}}
\newcommand{\cataloguescope}[1]{\paragraph{Mathematical question.}#1}
\begin{document}
% UnsolvedMath ID 8700061; source catalogue code AMR-086-0061
\setcounter{section}{8700060}
\section{Conjecture 4.19}\label{8700061}
{\small\sffamily UnsolvedMath ID 8700061 · Graph Theory}
\subsection{Definitions and mathematical statement}

\paragraph{Shared notation.}$\mathbb N=\{1,2,\ldots\}$, and $\mathbb Z,\mathbb Q,\mathbb R,\mathbb C$ denote the integers, rationals, reals and complex numbers. Any other conventions are specified locally.


Fix a real embedding $K\hookrightarrow\mathbb R$ of a number field and a simple abelian variety $A/K$ of dimension $g\ge1$. Let $A(\mathbb R)^0$ be the identity component of its real Lie group, and let $\ell=\operatorname{rank}_{\mathbb Z}A(K)$. Fix a symmetric ample line bundle on $A$, its Néron--Tate canonical height $\widehat h$, and a projective embedding $A\hookrightarrow\mathbb P^N$ from a sufficiently high tensor power of that line bundle. For real projective points $z=(z_0:\cdots:z_N)$ and $w=(w_0:\cdots:w_N)$ set
\[
\operatorname{dist}(z,w)=\frac{\max_{0\le i,j\le N}|z_iw_j-z_jw_i|}{\max_i|z_i|\max_j|w_j|}.
\]
\paragraph{Statement recorded in the supplied source.}
For every $\varepsilon>0$, does there exist $h_0>0$, depending on $A,K,\varepsilon$ (with the embedding and canonical height fixed), such that, for every real $h\ge h_0$ and every $\zeta\in A(\mathbb R)^0$, there is $\gamma\in A(K)$ satisfying
\[
\widehat h(\gamma)\le h,\qquad \operatorname{dist}(\zeta,\gamma)\le h^{-\ell/(2g)+\varepsilon}?
\] \sref{8700061}{2}

\subsection{Short English statement}
Can every real point in the identity component be uniformly approximated by rational points with bounded canonical height at this rank-dependent rate?
\subsection{Sources}
\begin{description}
\item[\hypertarget{source:8700061:1}{S1}] ULAM.AI and the UnsolvedMath Contributors, UnsolvedMath: A Curated Collection of Open Mathematics Problems, record 8700061 (AMR-086-0061). CC BY 4.0. \url{https://huggingface.co/datasets/ulamai/UnsolvedMath}.
\item[\hypertarget{source:8700061:2}{S2}] Michel Waldschmidt, \emph{Open Diophantine Problems}, arXiv:math/0312440v2 (24 January 2004), Conjecture 4.19 and projective-distance definition, p.43. \url{https://arxiv.org/abs/math/0312440v2}.
\end{description}
\label{end:8700061}
\end{document}
